\section{最適バックアップ設計}

投資 $I$ と追加保持世代数を固定すると，$\Delta$ に関する最適化は
\[
  G(\Delta;\lambda)=\frac{c_b}{\Delta}
  +\lambda\Qfull(\Delta;\alpha,\beta)
\]
の最小化に帰着する．ここで $\alpha=\alpha(I)$，
$\beta=\beta(I)$ であり，$\Qearly(I)$ と投資・保持費は
$\Delta$ に依存しない定数として除いている．

バックアップ間隔は数式上の正数全体ではなく，復旧時点目標や運用上の
サービス水準によって制限される．以下では
\[
  \mathcal{D}=[\Delta_{\min},\Delta_{\max}]
  \subset (t_m/n,\infty)
\]
を許容設計区間とする．

\begin{assumption}[内点設計条件]
\[
  \Lambda:=L-(d_1\tau_R+d_0)>0,\qquad
  \kappa_s>0,\qquad
  \Dstar:=\frac{\Lambda}{d_1\kappa_s n}
\]
とおく．許容設計区間は
$\Delta_{\min}<\Dstar<\Delta_{\max}$ を満たすとする．
$\kappa_s=0$ の場合，保持期間延長によるスキャン停止時間の限界費用が
消え，\eqref{eq:qfull-derivative} は内点零点を持たないため，本節の
閉形式内点解の対象外である．
\end{assumption}

\begin{theorem}[$\Qfull(\Delta)$ の閉形式微分と唯一の大域最小点]
$\Qfull(\Delta;\alpha,\beta)$ を \eqref{eq:qfull} で定義する．
仮定 5.1 のもとで
\begin{equation}
  \frac{\partial\Qfull}{\partial\Delta}
  =\eta n e^{-\beta n\Delta}f_{T_d}(n\Delta)
  \left(d_1\kappa_s n\Delta-\Lambda\right)
  \label{eq:qfull-derivative}
\end{equation}
が成立する．従って $\partial\Qfull/\partial\Delta=0$ の唯一解は
\begin{equation}
  \Dstar=\frac{L-(d_1\tau_R+d_0)}{d_1\kappa_s n}.
  \label{eq:dstar}
\end{equation}
また $\Delta<\Dstar$ では微分は負，$\Delta>\Dstar$ では正であり，
$\Dstar$ は $\mathcal{D}_n$ 上，したがって $\mathcal{D}$ 上でも
唯一の大域最小点である．さらに
\[
  \left.\frac{\partial^2\Qfull}{\partial\Delta^2}\right|_{\Dstar}
  =\eta d_1\kappa_s n^2 e^{-\beta n\Dstar}f_{T_d}(n\Dstar)>0 .
\]
\end{theorem}

\begin{proof}
$a(\Delta)=n\Delta$ とおく．$\Delta\in\mathcal{D}_n$ では
\[
  \Hbeta(\Delta;\alpha)
  =\int_{a(\Delta)}^\infty e^{-\beta t}f_{T_d}(t)\dd t,\qquad
  \Mbeta(\Delta;\alpha)
  =\int_{t_m}^{a(\Delta)} t e^{-\beta t}f_{T_d}(t)\dd t .
\]
Leibniz の積分則より
\[
  \frac{\partial\Hbeta}{\partial\Delta}
  =-n e^{-\beta n\Delta}f_{T_d}(n\Delta),\qquad
  \frac{\partial\Mbeta}{\partial\Delta}
  =n^2\Delta e^{-\beta n\Delta}f_{T_d}(n\Delta).
\]
\eqref{eq:qfull} を微分すると，$\Lbeta$ は $\Delta$ に依存しないので
\[
  \frac{\partial\Qfull}{\partial\Delta}
  =-\eta(d_1\tau_R+d_0)\frac{\partial\Hbeta}{\partial\Delta}
  +\eta d_1\kappa_s\frac{\partial\Mbeta}{\partial\Delta}
  +\eta L\frac{\partial\Hbeta}{\partial\Delta}.
\]
これを整理して \eqref{eq:qfull-derivative} を得る．
前因子 $\eta n e^{-\beta n\Delta}f_{T_d}(n\Delta)$ は正であるため，
符号は $d_1\kappa_s n\Delta-\Lambda$ のみで決まる．
よって零点は \eqref{eq:dstar} の一点に限られ，その前後で符号が
負から正へ変わる．
\end{proof}

\begin{remark}
$\Dstar$ は $\alpha,\lambda,\beta(I),\eta$ に依存しない．
これは早期検知失敗分岐の周辺点で，保持期間を少し延ばす便益
（全損から復旧可能へ移す効果）と追加スキャン停止時間の費用が，
同じ指数傾き因子と同じ症状検知確率 $\eta$ を持つためである．
ただし，微分の大きさと二階微分には
$\eta e^{-\beta n\Delta}f_{T_d}(n\Delta)$ が含まれるため，
有限 $\lambda$ での補正には尾指数，検知レート，症状検知確率が
影響する．
\end{remark}

\begin{theorem}[頻度漸近独立性]
$c_b>0$ とし，許容設計区間 $\mathcal{D}$ 上で
$G(\Delta;\lambda)=c_b/\Delta+\lambda\Qfull(\Delta;\alpha,\beta)$
の最小点を $\Delta^\ast(\lambda)$ と書く．このとき
\[
  \lim_{\lambda\to\infty}\Delta^\ast(\lambda)=\Dstar .
\]
\end{theorem}

\begin{proof}
$\mathcal{D}$ はコンパクトであり，$1/\Delta$ はその上で有界である．
従って
\[
  \frac{G(\Delta;\lambda)}{\lambda}
  =\Qfull(\Delta;\alpha,\beta)+\frac{c_b}{\lambda\Delta}
\]
は $\lambda\to\infty$ で $\Qfull(\Delta;\alpha,\beta)$ に
$\mathcal{D}$ 上一様収束する．定理 5.2 により
$\Qfull$ は $\mathcal{D}$ 上で唯一の最小点 $\Dstar$ を持つ．
一様収束と一意最小性から，任意の最小点列
$\Delta^\ast(\lambda)$ の極限点は $\Dstar$ に限られる．
コンパクト性により極限点は存在するので，全列が $\Dstar$ に収束する．
\end{proof}

\begin{corollary}[有限頻度での近似]
$Q''_\infty=\left.\partial^2\Qfull/\partial\Delta^2\right|_{\Dstar}$ とおく．
$\Delta^\ast(\lambda)$ が十分大きな $\lambda$ で内点にあるとき，
一階条件の Taylor 展開から
\[
  \Delta^\ast(\lambda)
  =\Dstar+\frac{c_b}{\lambda{\Dstar}^2 Q''_\infty}
  +O(\lambda^{-2})
\]
を得る．ここで
\[
  Q''_\infty
  =\eta d_1\kappa_s n^2 e^{-\beta n\Dstar}f_{T_d}(n\Dstar).
\]
\end{corollary}
